\chapter{La hacienda de un municipio pobre}

\begin{cautela}
Este capítulo mide lo que el municipio recaudaba. El capítulo \ref{cap:presencia} mide lo otro: qué puso el Estado provincial en el departamento en el mismo período ---qué escuelas, qué comisarías, qué oficinas y qué caminos--- y cuánto costó cada cosa. Los dos se leen mejor juntos, porque la renta de un municipio de tercera categoría sólo dice algo puesta al lado del gasto que otra jurisdicción hacía sobre el mismo territorio.
\end{cautela}


\label{cap:hacienda}

Un municipio de cuatro mil quinientos habitantes que en 2020 declaró no poder pagar los
sueldos de su personal administra hoy el suelo del área metropolitana. Este capítulo
reconstruye de dónde salen sus recursos, cuánto valen los inmuebles sobre los que cobra y
por qué lotear más no le aumenta el presupuesto.

\section{Cuánto recaudaba, en números}

Un decreto de 1938 permite ponerle cifra a esa hacienda que no alcanza.

\begin{ficha}{Decreto Nº 3161 del 10 de octubre de 1938}
Expediente 1117-M. Aprueba el balance de fondos del ejercicio \textbf{1937} de la Comisión
Municipal del Distrito de La Caldera y fija en \textbf{\$283,51} su contribución al fondo de
educación común, sobre la base del \textbf{diez por ciento afectable de la renta municipal}.

Firman Luis Patrón Costas y Víctor Cornejo Arias.
\end{ficha}

Si la contribución es el diez por ciento de la renta afectable, \textbf{la renta municipal de
1937 rondaba los dos mil ochocientos treinta y cinco pesos}. Es una inferencia aritmética y
como tal se la presenta: el decreto no publica el total.

Póngase esa cifra en su serie.

{\small
\begin{longtable}{@{}p{2.2cm}p{4.0cm}p{6.4cm}@{}}
\toprule
\textbf{Ejercicio} & \textbf{Renta municipal} & \textbf{Referencia contemporánea} \\
\midrule
\endhead
1908 & \$1.483,77 & Casi el ochenta por ciento venía de \textbf{vender en licitación} el derecho a cobrar los impuestos \\
1918 & \$3.387,50 \textit{(renta real)} & La calcula el decreto 1.660: cálculo de recursos de \$6.471,44 menos \$3.083,94 de saldo anterior «indebidamente computado», contra un mínimo de \$5.000 para ser municipalidad electiva (capítulo \ref{cap:siglo}) \\
1937 & \textbf{\$2.835} \textit{(inferido)} & Una casa quinta de diez hectáreas se remató en 1927 con base de \$6.000; la fundición de Mojotoro, en 1925, con base de \$222.144 \\
1946 & \$3.068 \textit{(recaudado)} & Recursos «efectivamente percibidos» que el decreto 5276-E de 1947 usa para repartir la coparticipación: al departamento le toca el 0,127 por ciento; a la capital, con \$1.298.124, el 54,116 (capítulo \ref{cap:siglo}) \\
\bottomrule
\end{longtable}}

En veintinueve años la renta municipal no llegó a duplicarse en términos nominales. \textbf{Y
la serie completa dice algo más}: entre 1908 y 1918 la renta más que se duplicó, y entre 1918 y
1946 no creció, sino que bajó. Las cuatro cifras no son del todo comparables ---una es la renta
real que calcula un decreto, otra una inferencia aritmética, otra lo recaudado que declara la
coparticipación---, pero ninguna de esas diferencias alcanza para explicar tres décadas sin
crecimiento nominal. Y seguía
siendo menos de la mitad de lo que valía una sola casa de fin de semana del departamento.

\textbf{Y la partida educativa seguía siendo la primera obligación.} En 1908 el veinte por
ciento adicional al Consejo General de Educación era la mayor erogación fija del ejercicio;
en 1932 el municipio arrastraba deuda por renta escolar atrasada y el decreto que aprobó su
presupuesto le ordenó tramitar el reconocimiento del saldo; en 1938 la contribución se fija
por decreto en el diez por ciento de la renta afectable.

\textbf{Treinta años, tres actos, la misma obligación y el mismo mecanismo}: el municipio
recauda poco, la Provincia le fija cuánto de eso va a educación, y el municipio lo debe.

\subsection*{Lo único que el archivo deja ver del régimen impositivo}

De las ordenanzas de impuestos del municipio no queda un solo artículo en el Boletín. Queda
una consulta, y su respuesta.

\begin{ficha}{Resolución Nº 434 del Ministerio de Gobierno --- Salta, 28 de abril de 1930 --- Expte.\ 637-M-1930}
«Vista la consulta elevada por la H.\ Comisión Municipal de La Caldera, sobre procedencia del
cobro de \textbf{impuestos de degolladura a los abastecedores de la peonada de la construcción
del camino nacional de Salta a Jujuy}»; y considerando que la Ordenanza General de Impuestos de
esa Municipalidad «no acuerda excepción alguna a favor de los referidos abastecedores, así como
tampoco existe Ley que la establezca», que «el impuesto de degolladura es de competencia
municipal» y que «los abastecedores que resisten su pago son comerciantes particulares»:

\textbf{Art.\ 1º.} «Declarar \textbf{procedente} el impuesto de degolladura aplicado por la
Municipalidad de Caldera, a los abastecedores de las cuadrillas de peones que trabajan en el
camino nacional de Salta a Jujuy, \textbf{dentro de la jurisdicción de aquélla}.»

Firma Rafael P.\ Sosa, Ministro de Gobierno. B.O.\ Nº 1.330 del 4 de julio de 1930.
\end{ficha}

Es el único acto de las ediciones leídas en que \textbf{el municipio de La Caldera cobra algo a un tercero
por lo que ocurre en su territorio, y el Estado provincial se lo confirma}. Y el tercero no es
cualquiera: son los proveedores de las cuadrillas que estaban construyendo, dentro del
departamento, el camino nacional a Jujuy ---la obra que el capítulo \ref{cap:infraestructura}
sigue durante todo el período.

Dos cosas se deducen de la resolución, y ninguna está dicha en ella. Que en 1930 existía una
\textbf{Ordenanza General de Impuestos} del municipio con articulado suficiente para que el
Ministerio verificara que no contemplaba excepciones ---de modo que el texto existía y alguien
lo leía---. Y que \textbf{hubo resistencia al pago}: la consulta nace de abastecedores que se
negaban.

\begin{cautela}
La resolución no transcribe ni un artículo de esa ordenanza, ni dice cuánto era el impuesto, ni
cuántos abastecedores ni por cuántas reses. El único contenido tributario municipal que este
relevamiento encontró en cuarenta años de Boletín es, entonces, el \textit{nombre} de un
impuesto y la confirmación de que se podía cobrar.
\end{cautela}

\textbf{Y un año antes, el Estado había autorizado a comprarle el terreno para la casa.} La Ley Nº 11088 de 1929 ---sancionada el 27 de septiembre y promulgada el 2 de octubre--- autoriza
la inversión de \textbf{\$500} para adquirir un terreno en el pueblo de La Caldera destinado a
la \textbf{Casa Municipal}, sin individualizar el terreno ni el vendedor; este libro no encontró el acto de compra. Quinientos pesos: lo mismo que seis años después se presupuestaría para el espigón
de emergencia para que el río no entrara a la calle principal.

\subsection*{Cinco actos seguidos, y ni una cifra}

Entre 1941 y 1945 el Boletín publica cinco actos consecutivos sobre la hacienda de La
Caldera. Los cinco aprueban. Ninguno publica un número.

{\small
\begin{longtable}{@{}p{3.5cm}p{2.1cm}p{5.0cm}p{2.0cm}@{}}
\toprule
\textbf{Acto} & \textbf{Fecha} & \textbf{Qué aprueba} & \textbf{Cifras} \\
\midrule
\endhead
Dto. 3072-G, expte. 1766-M/941 & 29/12/1941 & Balance del ejercicio económico 1940 de la Comisión Municipal & \textbf{ninguna} \\
Dto. 6148-G, expte. 190-M/943 & 14/06/1943 & Ordenanza General de Impuestos \textit{y} Ordenanza de Presupuesto de Gastos y Cálculo de Recursos, ejercicio 1943 & \textbf{ninguna} \\
Dto. 755-G, expte. 1549-M/943 & 30/09/1943 & Balance de la renta municipal del ejercicio 1942, elevado conforme el art.\ 89 de la Ley 68 & \textbf{ninguna} \\
Dto. 2857-G, expte. 6463-M/944 & 19/04/1944 & Cálculo de Recursos y Presupuesto de Gastos, ejercicio 1944 & \textbf{ninguna} \\
Dto. 5786-G, expte. 8266/944 & 12/01/1945 & Presupuesto de Gastos y Cálculo de Recursos, ejercicio 1945 & \textbf{ninguna} \\
\bottomrule
\end{longtable}}

Los cinco decretos tienen la misma estructura: visto, informes de Sanidad y del Consejo
General de Educación, dictamen del Fiscal de Gobierno, un artículo que aprueba, un artículo
que \textbf{devuelve el expediente a la comuna} y la fórmula de cierre. El de 1945 llega más
lejos que los otros cuatro: precisa dónde están los números ---«\textit{que corren a fojas 19 y
20 de las presentes actuaciones}»--- y no los transcribe.

\textbf{Se publica el acto; el contenido operativo queda en un expediente que no se publica.}
Es, ochenta años antes, la estructura que el capítulo \ref{cap:opacidad} describe para los
anexos de línea de ribera y para los contratos de promoción turística. Del régimen impositivo
municipal de La Caldera en 1943 y del monto de su presupuesto en 1944 y 1945, el Boletín no
permite saber absolutamente nada más que el hecho de su aprobación y la fecha.

\begin{cautela}
Las coordenadas de archivo, en cambio, quedan completas, y por eso el apéndice
\ref{cap:pedidos} las incorpora: expedientes 1766-M/941, 190-M/943, 1549-M/943, 6463-M/944 y
8266/944, todos devueltos por el propio articulado a la Municipalidad de La Caldera. El
presupuesto de 1945 está, si existe, en las fojas 19 y 20 del último.
\end{cautela}

\textbf{Y hay una vía por la que las cifras sí salieron a la luz: la deuda.} Lo que los cinco
decretos presupuestarios callan aparece, de a pedazos, en las leyes de crédito suplementario
que la Provincia sancionó en esos mismos años para cubrir renglones impagos del departamento
---entre ellos lo adeudado a la expendedora de rentas fiscales de La Caldera, en las Leyes 586
y 705, y el renglón de la Receptoría en la Ley 605---. \textbf{El municipio no publicaba su
presupuesto, pero la Provincia publicaba lo que había que pagarle a quien cobraba por él.} Es
la misma asimetría que el capítulo \ref{cap:opacidad} describe para el presente: lo que se
publica no es lo que se decide, sino lo que hay que saldar.

\textbf{Y hay un dato que corre en contra de la lectura fácil.} El presupuesto de 1943 se
aprueba en junio, con el ejercicio empezado; el de 1944, en abril; el de 1945, el \textbf{12 de
enero}. La demora se derrumba de cinco meses y medio a menos de dos semanas en tres años ---y
son los años de la intervención federal. La tutela presupuestaria no se ejerció siempre con la
misma lentitud, y quien quiera leer el retraso como medida de desidia tendrá que explicar por
qué el trámite más rápido de toda la serie ocurre bajo un gobierno de facto.

También queda fechado, entre esos mismos actos, un cambio institucional que este libro no
tenía. En junio de 1943 el decreto habla de la «\textit{H. Comisión Municipal del Distrito de
La Caldera}» y de su Presidente; en abril de 1944, del «\textit{Interventor de la Comuna de La
Caldera}». \textbf{El municipio fue intervenido en algún momento entre esas dos fechas}, y el
acto de intervención no está en las mil cuatrocientas noventa ediciones revisadas.

\subsection*{Y en 1980 el presupuesto seguía aprobándose desde afuera}

Cuarenta y ocho años después del decreto que aprobaba el presupuesto de 1932 con el ejercicio
a dos meses de cerrar, el mecanismo seguía funcionando igual.

El \textbf{Decreto Nº 1146} del 15 de agosto de 1980, del Ministerio de Gobierno, Justicia y
Educación, aprueba la \textbf{Ordenanza 028/79} de la Municipalidad de La Caldera, dictada el
27 de julio de \textbf{1979}. Total de recursos igual a total de erogaciones:
\textbf{\$258.178.081}. El mismo Boletín publica el equivalente para Vaqueros.

\textbf{El presupuesto de un ejercicio se aprueba trece meses después de sancionado y se
publica al año siguiente.} En 1932 la demora fue de diez meses; en 1980, de trece. La tutela
presupuestaria que el capítulo describe para 1908, 1922 y 1932 no es un rasgo de aquella
época: sigue vigente medio siglo después.

\textbf{Lo que no sigue una línea es la demora.} Puestos en serie ---1932, diez meses; 1943,
cinco y medio; 1944, tres y medio; 1945, dos semanas; 1980, trece meses---, los plazos suben y
bajan sin tendencia. Lo constante no es el retraso: es \textbf{quién aprueba}. En los cinco
casos, el presupuesto del municipio lo aprueba un decreto del Poder Ejecutivo provincial.

Y hay una comparación que conviene hacer con cuidado. Ese mismo año, la Provincia aprobó tres
obras de encauzamiento en el departamento por quinientos veintisiete millones de pesos ---
\textbf{nominalmente, el doble del presupuesto anual del municipio sancionado un año antes}.

\begin{cautela}
Las dos cifras están en la misma moneda, pero no en el mismo momento: el presupuesto se sancionó en julio de 1979 y las obras se aprobaron en agosto y septiembre de 1980, con una inflación anual que en esos años superó el cien por ciento. En términos reales la relación puede estar cerca de uno a uno, y la comparación debe leerse con esa reserva. Pero no debe leerse como que el municipio «recibió» ese dinero: \textbf{las obras
las contrató, ejecutó y pagó la Provincia}, y el municipio no interviene en ninguna de las
tres resoluciones.

Eso es precisamente lo que la comparación muestra. Sobre el territorio de un municipio cuyo
presupuesto anual era de doscientos cincuenta y ocho millones, se decidieron en tres semanas
intervenciones por más del doble de esa cifra, sin que el municipio fuera parte.
\end{cautela}

\subsection*{Y en 1940 también debía sanidad}

La Resolución Nº 2025 del 29 de mayo de 1940 publica la nómina de municipios deudores de la
Dirección Provincial de Sanidad al 31 de marzo de ese año. En el renglón quince figura
\textbf{La Caldera, con \$951,95}.

Puesta contra la renta municipal inferida para 1937 --- unos dos mil ochocientos treinta y
cinco pesos ---, esa deuda equivale a \textbf{alrededor de un tercio de todo lo que el
municipio recaudaba en un año}.

Y se suma a la anterior. En 1932 el municipio arrastraba deuda por renta escolar atrasada; en
1938 su aporte educativo se fijó por decreto en el diez por ciento de la renta afectable; en
1940 debe casi mil pesos a Sanidad. \textbf{Las dos únicas obligaciones que el archivo
documenta --- educación y salud --- las debía simultáneamente.}

\subsection*{Las rendiciones aparecieron}

El apéndice \ref{cap:pedidos} de este libro reclamaba la rendición de cuentas del ejercicio 1933, que en su
momento se había desglosado del expediente «para ser considerada por separado» y no volvió a
aparecer. Volvió: un decreto del 18 de octubre de 1934, publicado en junio de 1935, aprueba
\textbf{las rendiciones de los ejercicios 1932 y 1933} de la Comisión Municipal de La
Caldera, en los expedientes 1741-M, 1258-M y 1256-M, conforme el artículo 89 de la Ley 68.

Se aprueban \textbf{«con las modificaciones señaladas por la Contaduría del Consejo General de
Educación»}. Firman A.\ Aráoz y Alberto B.\ Rovaletti.

Dos cosas.

\textbf{Quien observa las cuentas del municipio es la Contaduría del Consejo de Educación},
no la Contaduría General de la Provincia. El organismo ante el cual el municipio arrastraba
deuda es el mismo que audita su rendición. No hay nada irregular en eso --- la contribución
educativa era la principal obligación del municipio y el organismo acreedor tenía interés
legítimo en revisarla ---, pero completa el cuadro de dependencia que este capítulo describe.

\textbf{Y la aprobación llega dos años tarde.} Las cuentas de 1932 se aprueban en octubre de
1934 y se publican en junio de 1935. Sumado a que el presupuesto de 1932 se había aprobado
en octubre de ese mismo año --- con el ejercicio a dos meses de cerrar ---, el ciclo completo
es el de una administración que autoriza tarde y controla después.

\textbf{Y las rendiciones siguieron apareciendo, siempre igual.} El barrido continuo de
1935--1945 agrega las de los ejercicios \textbf{1935} y \textbf{1937} ---esta última devuelta
por el Decreto 2155 y aprobada después--- y la de \textbf{1940}, que es a la vez el primero de los cinco actos de la sección anterior. Sumadas a las de 1932 y 1933
y a los cinco actos presupuestarios de la sección anterior, son \textbf{nueve actos sobre la
hacienda de La Caldera entre 1932 y 1945} ---contando por separado las rendiciones de 1932 y 1933, que aprobó un solo decreto---, todos publicados y sólo uno con una cifra: el Decreto 3161/1938, que fija la contribución educativa sin publicar la renta. Lo que el
Boletín deja ver de las finanzas del municipio en trece años es la lista de los expedientes
que las contienen.

\section{El municipio como contratista, en 1933}

Este capítulo documenta que en 2014 la Dirección de Vialidad de Salta contrató directamente
a la Municipalidad de La Caldera para construir badenes, desmalezar banquinas y levantar
gaviones en ocho rutas provinciales, y que ese contrato fue parte de un programa que alcanzó
a nueve municipios el mismo día. La figura tiene ochenta y un años más de los que parece.

\begin{ficha}{Acta Nº 18 del Directorio de Vialidad de Salta, 20 de junio de 1933}
Punto 14: «Resultando muy peligroso el tráfico de carros y jardineras por el \textbf{camino
en cornisa de Salta a Jujuy}, debido al poco ancho de su calzada y a fin de subsanar el
referido inconveniente que perjudica enormemente el transporte de productos de la zona que
atraviesa dicho camino, este Directorio resuelve invertir hasta la suma de \textbf{quinientos
pesos} en la \textbf{limpieza del camino de La Caldera a Tres Cruces por el bajo},
\textbf{encomendando la Dirección de éste trabajo al Presidente de la Municipalidad de La
Caldera D.\ Augusto Regis\index{Regis, Augusto}} quien presentará en su oportunidad toda la documentación
referente a la inversión de los fondos acordados, a cuyo efecto se le remiten las planillas
necesarias.»

Presidencia del Directorio: ingeniero José Alfonso Peralta. Vocales: Domingo Patrón Costas,
Arturo Michel, Sergio López Campo. El acta se publica como anexo del decreto que la eleva a
aprobación del Poder Ejecutivo.
\end{ficha}

El monto es menor. La figura no lo es, y conviene descomponerla.

\textbf{El problema es de la cornisa, y la solución no es la cornisa.} El camino nacional de
Salta a Jujuy es peligroso por angosto; en lugar de ensancharlo, se acondiciona una vía
alternativa por el bajo para descargarlo. Es la misma lógica que el capítulo \ref{cap:defensas} documenta en
1910, cuando la Provincia interviene sobre un camino nacional porque «el Gobierno de la
Nación no tiene votados fondos especiales a este objeto».

\textbf{Y la ejecución se delega en el municipio.} No se hace por administración ni por
contrato con una empresa: se le encomienda al presidente municipal, con planillas y rendición
posterior. En 1933 como en 2014, el municipio es el brazo caminero de una jurisdicción mayor
sobre un camino que no es suyo.

Hay una diferencia, y juega a favor de 1933. El acta nombra a la persona, fija el tope, exige
documentación y remite las planillas. Es decir: hay control de rendición explícito en el
propio acto que asigna. Los contratos de 2014 informan el resultado de una contratación
directa por libre negociación y no publican ni el certificado de obra ni el acta de
recepción, que este libro reclama en el apéndice \ref{cap:pedidos}.

\section{Diez años sin impuestos por plantar olivos}

Los capítulos \ref{cap:poblacion} y \ref{cap:tierrafiscal} documentan que en 2014 y 2015 la Provincia eximió de tributos provinciales a
dos hoteles del pueblo, y que en 2015 cedió cuarenta y una hectáreas fiscales en comodato
gratuito para una cancha de golf. La exención como instrumento de política económica sobre
este departamento tiene ochenta años más.

\begin{ficha}{Ley provincial Nº 145 (numeración original), sancionada el 2 de octubre de 1934}
\textbf{Art.\ 1º:} «Formando parte integrante de la \textbf{región económica del olivo}, los
Departamentos de La Candelaria, Rosario de la Frontera, Metán, Cafayate, Guachipas, San
Carlos, La Viña, Chicoana, Rosario de Lerma, Cerrillos, Capital y \textbf{Caldera}, de acuerdo
con lo establecido por la Ley Nacional Nº 11643, declárase a la Provincia \ldots\ acogida a
los beneficios de la citada Ley.»\par\textit{El 145 es un número de la serie «original» de los años treinta. El Digesto de la Cámara de Diputados reserva «Ley 145» para una ley de 1868 y registra las de esta serie con otra numeración ---la Ley 1340, por ejemplo, como «original 59»---. En la tabla oficial de leyes con doble numeración figura como \textbf{Ley 1423 (orig.\ 145)}, publicada el 12 de octubre de 1934; con ese número hay que pedirla.}

\textbf{Art.\ 2º:} «Quedan \textbf{eximidas por el término de diez años, de todo impuesto
provincial ó municipal, las tierras en que se planten olivos}.»
\end{ficha}

Tres cosas.

\textbf{Hubo un intento de asignarle una vocación productiva al departamento}, por ley y con
beneficio nacional. La Caldera fue incluida en la región económica del olivo junto a otros
once departamentos. El capítulo \ref{cap:tierra} documenta qué quedó de la agricultura departamental
ochenta y cuatro años después: una pérdida del ochenta y ocho por ciento de la superficie
agropecuaria mientras la provincia ganaba. El olivo no está entre lo que se cultiva hoy.

\textbf{La Provincia eximió de impuestos municipales sin intervención del municipio.} El
artículo 2º libera esas tierras de todo tributo provincial \textbf{o municipal} por diez
años. Es decir: una ley provincial resigna la recaudación de un municipio que no la votó, en
un momento en que ese mismo municipio arrastraba deuda escolar y gastaba en sueldos más de
lo que la ley le permitía.

No hay nada excepcional en eso --- las provincias legislan sobre tributos municipales con
frecuencia --- y este libro no lo objeta. Lo registra porque completa el cuadro de la
autonomía: el municipio no decidía su categoría, no aprobaba su presupuesto, no fijaba su
planta, y tampoco controlaba las exenciones sobre su propia base imponible.

\textbf{Y la exención es antigua como técnica.} En 1934 se eximía por plantar olivos; en
2014 y 2015, por ampliar hoteles. El instrumento es el mismo y el destinatario cambió: de
quien produce en la tierra a quien aloja sobre ella. Es, en una sola línea de continuidad
fiscal, el arco que este libro describe.

\section{El presupuesto llegaba tarde y con deuda}

El capítulo \ref{cap:siglo} documenta que en 1922 el Poder Ejecutivo provincial aprobaba por decreto el
presupuesto que la comisión municipal sancionaba. Diez años después el mecanismo seguía
igual, y un decreto permite ver qué significaba en la práctica.

\begin{ficha}{Decreto Nº 15.445 del 28 de octubre de 1932}
Expediente 1392-M, dictado en acuerdo de ministros. Aprueba la Ordenanza General de
Impuestos, el Presupuesto de Gastos y el Cálculo de Recursos elevados por la
\textbf{Municipalidad de 3ª Categoría de La Caldera}, «proyectados para regir durante el
ejercicio económico del presente año 1932».

Previo dictamen fiscal del 12 de agosto e informe de la Contaduría del Consejo General de
Educación del 25 de octubre. No se formula observación alguna.

\textbf{Dos cargas impuestas al municipio:} elevar al Poder Ejecutivo la rendición de cuentas
del ejercicio 1931, conforme el artículo 184 de la Constitución; e \textbf{iniciar ante el
Consejo de Educación la tramitación para establecer el saldo deudor por renta escolar
atrasada}, conforme la ley de emisión de obligaciones de la Provincia.

Firman Aráoz, Rovaletti y García Pinto (h).
\end{ficha}

Dos detalles de ese acto describen la posición del municipio mejor que cualquier
generalización.

\textbf{El presupuesto se aprobó el 28 de octubre para regir durante 1932.} Es decir, con el
ejercicio a dos meses de terminar. Un municipio que recibe la autorización para gastar cuando
ya gastó no planifica: rinde. Y el decreto que lo aprueba deja constancia, en el mismo acto,
de que todavía debía la rendición del año anterior.

\textbf{Y arrastraba deuda escolar.} El municipio no había girado la contribución que le
correspondía al Consejo General de Educación, y el decreto le ordena tramitar el
reconocimiento de ese saldo. La partida existía desde 1908 --- el cuadro de rentas de aquel
año consigna «al Consejo General de Educación, el 20 por ciento adicional, \$200», la mayor
erogación fija del ejercicio --- y veinticuatro años después seguía impaga.

Y al año siguiente el mecanismo se repite con un agregado que importa. El \textbf{Decreto
Nº 16.798}, del 2 de octubre de 1933 y publicado tres meses más tarde, aprueba el
presupuesto municipal del ejercicio 1933 pero deja constancia de una observación fiscal: el
presupuesto de gastos \textbf{contraría el artículo 79 de la Ley Nº 68} ---la Ley Orgánica de Municipalidades, que el Digesto registra hoy como Ley 1349 (apéndice \ref{cap:normativa})---, que prohíbe invertir
más del veinticinco por ciento de la renta en sueldos de empleados administrativos.

Se aprueba igual, «con la modificación señalada en el dictamen fiscal». Y la rendición de
cuentas del municipio se desglosa del expediente «para ser considerada por separado»: ese
trámite separado terminó en el decreto del 18 de octubre de 1934 que se analiza más arriba.

De modo que en 1933 el municipio de La Caldera gastaba en sueldos administrativos más de lo
que la ley permitía, la Provincia lo advirtió y lo aprobó de todas formas. Ochenta y siete
años después, en julio de 2020, ese mismo municipio declaró por escrito que no podía pagar
los sueldos de su personal.

\textbf{No es la misma situación} --- una es exceso de planta sobre renta, la otra es falta
de caja --- pero las dos describen el mismo desajuste estructural: la estructura
administrativa mínima de este municipio nunca entró en su recaudación.

Ninguna de estas cosas es un escándalo administrativo. Son el retrato ordinario de una
hacienda que no alcanza. Pero conviene retenerlas cuando este capítulo llegue a 2020 y
encuentre a ese mismo municipio declarando por escrito que no puede pagar los sueldos de su
personal: la incapacidad de sostener su propia estructura no aparece con el crecimiento
demográfico. Es anterior a él en casi un siglo.


\section{Julio de 2020: el municipio no puede pagar sueldos}

En julio de 2020, empleados municipales de La Caldera cortaron la Ruta Nacional 9 en apoyo de siete trabajadoras cuyos contratos no fueron renovados. Según \textit{Salta 4400}, el intendente informó que los contratos habían vencido en abril, que continuaron trabajando en mayo y junio para no dejarlas sin empleo, y que la situación económica no daba para más: sostuvo que no tenía dinero para recontratarlas y que se encontraba con los números en rojo para pagar los sueldos del resto del personal, porque la coparticipación ya no alcanzaba.

Las trabajadoras atribuyeron la medida a razones políticas y formularon otras imputaciones sobre la composición de las contrataciones. Este libro no las verifica ni las reproduce: retiene del episodio el reconocimiento público, por el propio Ejecutivo municipal, de una situación de insolvencia corriente.

El dato importa por su fecha. En julio de 2020 el municipio declara que la coparticipación no le alcanza para pagar salarios. Cuatro años antes, el municipio vecino lo había dejado escrito en dos decretos: en octubre de 2016 el intendente de Vaqueros promulgó las ordenanzas que mandaban construir una pasarela en el acceso al barrio Villa Sara y una ciclovía en Lesser (Ordenanzas 502/16 y 503/16, Decretos 156/16 y 157/16) aclarando que su ejecución no cabía en el presupuesto de 2017 «dada la magnitud» de las obras y por «la escasa coparticipación que se recibe por un municipio pequeño como el nuestro», y la subordinó a fondos provinciales o nacionales. Cinco años después, el puente peatonal del ingreso a Villa Sara figuraba entre los proyectos de la senadora por el departamento (capítulo \ref{cap:politica}). La Caldera, por su parte, cinco años antes había aprobado un Plan de Desarrollo Urbano Ambiental y estaba por reglamentarlo con un Código de 152 artículos que crea un Órgano Técnico de Aplicación transdisciplinario, un procedimiento de evaluación de impacto ambiental de dieciséis artículos y un régimen de urbanización de quince pasos con certificados sucesivos, control de avance y garantías de obra.

Ese instrumento presupone una capacidad estatal que el municipio, por su propia declaración, no tiene. No se trata de desidia ni de captura: \textbf{el Código de La Caldera fue diseñado para un municipio que no existe}. La Provincia y el Consejo Federal de Inversiones financiaron y asistieron técnicamente la redacción de la norma; nadie financió la estructura que debía aplicarla.

Ésa es la segunda respuesta --- y la más económica --- a la pregunta que abrió el capítulo \ref{cap:cot}. Es también la que menos sirve para una denuncia, porque no tiene responsable individual.

\subsection*{Cuánto vale un lote baldío, según el Estado}

La discusión sobre la base imponible municipal tiene, en un edicto de remate de 2015, una cifra exacta.

\begin{ficha}{Remate judicial --- B.O.\ Nº 19.651 a 19.653, octubre de 2015}
En los autos ``Huaranca Carvajal, Rocío del Valle vs.\ Goitia, Ignacio --- Ejecución de sentencia'', expediente 434.385/13 del Juzgado Civil y Comercial de 6ª Nominación, se rematan:

el cincuenta por ciento indiviso del \textbf{catastro 1411}, sección C, manzana 9, parcela 11 del departamento La Caldera, ``con la base de las dos terceras partes del valor fiscal, \textbf{o sea \$152,88}''; y el cincuenta por ciento indiviso del \textbf{catastro 1412}, parcela 12, con base de \textbf{\$185,48}.

``Ambos inmuebles se encuentran \textbf{baldíos y deshabitados}. Servicios: \textbf{sin conexión de agua corriente ni energía eléctrica}. Estado de ocupación: desocupados.''

Costo de publicación del edicto: \textbf{\$420}.
\end{ficha}

Hágase la cuenta, porque es de una elocuencia que ningún argumento mejora.

La base de remate es dos tercios del valor fiscal. De modo que el valor fiscal completo de esas mitades indivisas era de unos \$229 y \$278; el de los inmuebles enteros, del orden de \$458 y \$556. \textbf{Los dos lotes juntos estaban valuados fiscalmente, en 2015, en poco más de mil pesos.}

\textbf{El aviso que anunció su remate costó \$420.} Publicar la venta de dos lotes urbanos costó, en el Boletín Oficial, cuatro décimos de lo que el Estado dice que valen.

Ahí está, con número, lo que este capítulo viene sosteniendo. El impuesto inmobiliario y las tasas municipales se calculan sobre esa valuación. Un municipio cuya superficie urbana se multiplica en lotes valuados a cuatrocientos o quinientos pesos no aumenta su recaudación por lotear: \textbf{aumenta su superficie a servir sin aumentar su base imponible}. Cada lote nuevo suma obligación de servicio y aporta un tributo calculado sobre una cifra simbólica.

Y la descripción judicial confirma por otra vía lo que el censo mide: baldíos, deshabitados, sin agua ni luz. No es una excepción encontrada a propósito --- es la descripción rutinaria de un martillero sobre dos parcelas cualesquiera del departamento.

\begin{cautela}
Una advertencia sobre el alcance de este dato, porque conviene no estirarlo.

Son \textbf{dos parcelas}, no una muestra. Su valuación puede estar desactualizada respecto del padrón general, puede corresponder a una zona particularmente marginal del ejido, o puede reflejar un revalúo pendiente que después se corrigió. Una valuación fiscal baja tampoco equivale a un valor de mercado bajo: precisamente el desfasaje entre ambos es lo que este libro sospecha.

Lo que la cifra permite afirmar es más modesto y sigue siendo suficiente: que en 2015 existían en el departamento lotes urbanos cuya valuación fiscal era \textbf{de tres cifras}, y que sobre esas valuaciones se liquidan los tributos que financian al municipio. El padrón inmobiliario completo diría si son la regla o la excepción, y es uno de los pedidos que este capítulo formula.
\end{cautela}


\subsection*{Y en 2007 el municipio pagaba para saber qué había construido}

La base imponible no es sólo una cifra: es una información que el municipio no produce.

\begin{ficha}{Decreto Nº 1639/07 --- B.O.\ Nº 17.645, 19 de junio de 2007}
Aprueba las Actas Acuerdo entre la Provincia, por el Ministerio de Hacienda y Obras Públicas, y los municipios de Vaqueros (22/03/2007, intendente Miguel Ángel Alemán), La Caldera (17/04/2007, intendente Héctor Miguel Calabró) y El Galpón, para que la Dirección General de Inmuebles les preste el \textbf{Servicio de Información Catastral}: una base georreferenciada, con relevamiento satelital y aerofotogramétrico, que detecta las mejoras construidas y no declaradas y las valúa.

Obligaciones del municipio: informar la base catastral y tributaria existente, incorporar las superficies edificadas que le informe la Provincia y comunicarle las mejoras que conozca. \textbf{Precio: el 30 por ciento de lo que el municipio recaude de más por las mejoras detectadas}; en el acta de La Caldera, el 15 por ciento hasta que el municipio se incorpore al Plan de Fortalecimiento Tributario Municipal. Plazo anual, prorrogable tácitamente hasta que se dicte una ley que regule el servicio.

Anexo I de las dos actas, con el mismo texto: el plan de entrega de la información se ordenará por la sectorización del plano de Inmuebles, «considerando que \textbf{La Municipalidad no cuenta con Código de Planeamiento Urbano actual}».
\end{ficha}

\begin{cautela}
Tres cosas. La primera es de fecha: en 2007, según actas firmadas por sus propios intendentes, ni La Caldera ni Vaqueros tenían código de planeamiento. Los dos aprobaron su plan en noviembre de 2015 y su código en 2016, el de Vaqueros el 26 de diciembre (capítulos \ref{cap:cot} y \ref{cap:vaqueros}). El ciclo de loteos que este libro documenta empezó antes que cualquiera de ellos.

La segunda es de dependencia. El impuesto inmobiliario urbano es un recurso propio del municipio, pero la valuación la hace la Provincia, y el municipio pagaba con una parte de su propia recaudación para que la Provincia le dijera qué se había construido en su territorio. Es la forma tributaria de la misma asimetría que este capítulo describe: la función es municipal; el instrumento, provincial.

La tercera es una pregunta. Las actas preveían una Unidad de Seguimiento y Control con dos funcionarios por parte, y una ley que regularía el servicio para todos los municipios. Esa ley no aparece: la búsqueda de leyes provinciales por «información catastral» devuelve sólo la Ley de Catastro de 1948 y otras normas ajenas al tema. Si el proyecto no se sancionó, el régimen quedó librado a convenios anuales prorrogables por silencio. Si el servicio se prestó, cuántas mejoras detectó en el departamento y cuánto pagó el municipio por ellas es lo que dirían los informes de esa unidad. \pendiente{informes de la Unidad de Seguimiento y Control del Servicio de Información Catastral para La Caldera y Vaqueros, y destino del proyecto de ley de 2007 en la Legislatura}
\end{cautela}


\subsection*{El municipio como contratista de la provincia}

\begin{ficha}{Contratación directa por libre negociación --- B.O.\ Nº 19.405, 15 de octubre de 2014}
La Dirección de Vialidad de Salta informa el resultado de la contratación directa por libre negociación \textbf{con la Municipalidad de La Caldera}, expediente 33-192.402/2014-0, encuadrada en los arts.\ 13 y 24 de la Ley 6.838. Presupuesto oficial \textbf{\$307.877}. Contrato aprobado por Resolución 1616/2014, desde el 1 de septiembre de 2014 y por \textbf{cuatro meses}.

\textbf{Trabajos:} construcción de \textbf{badenes y alcantarillas}, \textbf{desmontes de zona de camino}, desmalezado de banquinas, limpieza de alcantarillas, desembanque y limpieza de obras de arte, construcción de \textbf{gaviones de piedra, colchonetas y muros de piedra}, cuadrilla y camión.

\textbf{Alcance:} Rutas Provinciales 11, 28, 107-s, 114-s, 115-s, 124-s, 125-s, 150-s, Acceso Norte, Circunvalación Oeste, Centro de Convenciones y caminos vecinales.

En la misma edición se publican contrataciones idénticas con las municipalidades de El Jardín, La Viña, La Candelaria, Apolinario Saravia, Chicoana, El Quebrachal, San Lorenzo y El Bordo. Y la modalidad se repite al año siguiente: por Resolución 0207/2015 la misma Dirección contrató nuevamente al municipio, desde el 5 de enero de 2015 y por cuatro meses, por \$264.815, con el mismo repertorio de tareas más un camión volcador de seis metros cúbicos ``modelo no inferior al año 1986''.

Y no es la única modalidad. Por \textbf{Resolución SOP Nº 729 del 12 de septiembre de 2014} la Provincia adjudicó a la misma Municipalidad de La Caldera, con encuadre en el art.\ 13 inc.\ a de la Ley 6838, la ejecución de la obra ``Construcción de dos aulas y galería'' en el \textbf{Colegio Nº 5045 `Senado Provincial'} de La Caldera, por \textbf{\$625.663,90} y ciento cincuenta días de plazo, con convenio y crédito legal a favor del municipio.
\end{ficha}

\begin{cautela}
Tres lecturas.

\textbf{Es un programa, no una excepción.} Nueve municipios contratados el mismo día por montos casi iguales --- el de La Caldera coincide al peso con el de Chicoana --- indican una política provincial de tercerizar el mantenimiento vial en los municipios, no un acuerdo singular. Este libro no lo objeta: es una forma razonable de sostener caminos secundarios con estructuras locales.

\textbf{Y da una medida de escala.} Entre el mantenimiento vial y la obra escolar, la Provincia canalizó por el municipio más de novecientos treinta mil pesos en cuatro meses de 2014. Para un municipio de cuatro mil quinientos habitantes es un flujo significativo, y no proviene de la coparticipación ni de sus tasas: proviene de \textbf{prestar servicios y ejecutar obra por cuenta del Estado provincial}, con su propia cuadrilla y su propio camión. El municipio no sólo administra: es contratista.

\textbf{Lo tercero es lo que conecta.} Entre los trabajos contratados figuran ``construcción de badenes'' y ``desmontes de zona de camino'', sobre rutas y caminos vecinales del departamento. Son, término por término, las mismas tareas que el municipio informa estar ejecutando con maquinaria pesada en febrero de 2026, diez días antes de la audiencia judicial.

De modo que la capacidad municipal para construir badenes y desmontar zonas de camino \textbf{no es reciente ni improvisada}: existe desde al menos 2014, está financiada por Vialidad provincial y se ejerce por contrato. Eso no dice nada sobre la legitimidad de lo hecho en 2026 --- que este libro no juzga ---, pero descarta una explicación posible: que las obras de 2026 fueran una respuesta improvisada de un municipio sin medios.
\end{cautela}


\subsection*{Y en 2017 seguía siendo contratista de Vialidad}

Este capítulo documenta que en 1933 la Dirección de Vialidad le encomendó al
presidente municipal la limpieza de un camino, y el contrato de 2014
por badenes y gaviones en ocho rutas provinciales. Hay un tercer punto en la serie.

En octubre de 2017, la Dirección de Vialidad de Salta publica \textbf{dos contrataciones
directas con la Municipalidad de La Caldera}, ambas encuadradas en el artículo 13 de la Ley
6.838, por \textbf{\$606.078} y \textbf{\$493.878}: mantenimiento vial en cinco rutas
provinciales, caminos vecinales de la zona y el acceso al Centro de Convenciones, más
provisión de camión y de maquinaria.

En los dos casos la \textbf{cotización es idéntica al presupuesto oficial}, hasta el centavo.
Los plazos son de cuatro y tres meses. Firma el ingeniero Gerardo R.\ Villalba.

\textbf{La serie queda así: 1933, 2014 y 2017.} Ochenta y cuatro años en los que el municipio
es contratado directamente por el organismo vial de otra jurisdicción para mantener caminos
que no son suyos.

Y conviene retener el dato de escala: un millón cien mil pesos en dos contratos, para un
municipio de tercera categoría que tres años después declararía por escrito no poder pagar los
sueldos de su personal. \textbf{La obra vial para terceros no es un ingreso marginal de esta
hacienda: es una de sus fuentes.}

\subsection*{Y en 2019 la ruta 11 la mantenía otro municipio}

\begin{ficha}{Dirección de Vialidad de Salta --- contrataciones con la Municipalidad de El Bordo, 2019}
Tres contrataciones sucesivas para el mantenimiento de las Rutas Provinciales 10-P, \textbf{11-P, tramo «La Calderilla -- El Desmonte»}, y 12-P: desmalezamiento de zona de camino, bacheo con material bituminoso y provisión de camión.
\begin{itemize}[leftmargin=*,nosep]
\item Contratación directa por libre negociación, expte.\ 33-17063/2019, Res.\ Vialidad 251/19: \$776.282,34, tres meses desde el 1/2/2019 (B.O.\ 12/02/2019).
\item Contratación abreviada, expte.\ 33-90807/2019, Res.\ 875/19: \$1.211.028,08, cuatro meses desde el 2/5/2019 (B.O.\ 14/06/2019).
\item Contratación abreviada, expte.\ 33-189497/2019, Res.\ 1214/2019: \$756.887,74, dos meses desde el 1/7/2019 (B.O.\ 16/08/2019).
\end{itemize}
Las dos últimas se encuadran en el artículo 15, inciso a), de la Ley 8072 y en el artículo 40 del Decreto 1319/18.
\end{ficha}

\begin{cautela}
La Ruta Provincial 11 figuraba en 2014 entre las que Vialidad le encomendaba a la Municipalidad de La Caldera, y es la misma que el plano de Chacras del Gallinato desmembra en 2015 (capítulo \ref{cap:loteo}). En 2019, el tramo que arranca en La Calderilla lo mantiene por contrato un municipio de otro departamento, El Bordo, por unos 2,7 millones de pesos en siete meses. Este libro no sabe si La Caldera dejó de ser contratista para esa ruta, si el tramo contratado a El Bordo cruza el departamento o empieza en su límite, ni si hubo en 2019 contratos con La Caldera por otras rutas; los títulos de los avisos de ese año no los muestran. Lo que sí muestran es que el mantenimiento de un camino que nace en el departamento se decide y se paga sin pasar por su municipio.

Una nota de encuadre: el mismo artículo 15 de la Ley 8072 que la Provincia invocó en 2024 para contratar la consultoría del plan de manejo (inciso e, capítulo \ref{cap:amparo}) aparece aquí con otro inciso, el a), para contratar con un municipio. \pendiente{texto del artículo 15 de la Ley 8072 y contratos de Vialidad con la Municipalidad de La Caldera posteriores a 2017}
\end{cautela}

\subsection*{Y una tercera fuente de ingresos: las licencias}

\begin{ficha}{Decreto Nº 1433 del 29 de abril de 2015 --- B.O.\ Nº 19.531}
Expediente 342-39.526/14. Aprueba el Convenio Específico de Distribución del \textbf{Fondo de Cooperación Técnica y Financiera} entre la Provincia --- representada por el Ministro de Seguridad --- y la \textbf{Municipalidad de La Caldera}, representada por su intendente \textbf{Luis Gerardo Mendaña\index{Mendaña, Luis Gerardo}}.

Por ese convenio la Provincia \textbf{cede y autoriza al municipio a percibir directamente el cien por ciento} de los montos que le corresponderían por la venta de formularios del \textbf{Certificado Nacional de Antecedentes de Tránsito}, correspondientes al \textbf{Centro de Emisión de Licencias} de su jurisdicción.

Antecedente: el convenio del 2 de junio de 2010 entre la Agencia Nacional de Seguridad Vial, la Provincia y el municipio, aprobado por Decreto 4257/11, que certificó el Centro de Emisión de Licencias municipal.

\textbf{El convenio} ---el Anexo, vinculado al decreto en la edición web como copia escaneada--- \textbf{se firmó el 10 de marzo de 2014}, más de un año antes de su aprobación. Del valor bruto de cada formulario, el treinta y seis por ciento corresponde a la Provincia y al municipio; la Provincia cede su parte, de modo que el municipio percibe \textbf{todo ese treinta y seis por ciento}, \textbf{con efecto desde la implementación del sistema} y con \textbf{destino exclusivo a seguridad vial} en su jurisdicción. Su vigencia es la del convenio de 2010, y cualquiera de las partes puede rescindirlo con tres meses de aviso.
\end{ficha}

Un municipio de cuatro mil quinientos habitantes tiene \textbf{Centro de Emisión de Licencias de conducir} desde al menos 2010, conectado a los registros nacionales, y se queda con la porción del formulario que corresponde a la jurisdicción, con destino obligado a seguridad vial.

Sumado a lo anterior, el cuadro de recursos del municipio queda así: coparticipación y tasas propias; contratos de mantenimiento vial y obra escolar con la Provincia; y emisión de licencias con retención íntegra del arancel. \textbf{Al menos dos de esas tres fuentes no dependen de su base imponible sino de convenios con el Estado provincial}, y la de licencias, además, sólo puede gastarse en seguridad vial.

Eso matiza --- sin desmentirla --- la afirmación de que la base fiscal del municipio crece en superficie loteada y no en recaudación. Es cierto respecto de las tasas, pero el municipio compensa esa debilidad con ingresos convenidos. Lo que no cambia es la conclusión de fondo: son recursos que no derivan del suelo que administra, y por lo tanto \textbf{más loteo no significa más presupuesto}.


